\documentclass[9pt]{extarticle}

\usepackage[utf8]{inputenc}
\usepackage[a4paper, margin=1.2cm]{geometry}
\usepackage[colorlinks=true, urlcolor=blue]{hyperref}
\usepackage{enumitem}
\usepackage{changepage}

\pagestyle{empty}

\setlist[itemize]{label=\textbullet, noitemsep, topsep=0pt}
\setlength{\parindent}{0pt}

\begin{document}

\begin{center}
	{\huge\bfseries Paul Parisot}
	\\[4pt]
	{\Large Ingénieur Logiciel - C++, Java \& Finance}
	\\
\end{center}

\par\vspace{-5pt}
\rule{\linewidth}{0.4pt}
\par\vspace{-5pt}

\begin{center}
	Github: \href{https://github.com/paulogarithm}{paulogarithm}
	\\
	\href{mailto:paul.parisot@epitech.eu}{paul.parisot@epitech.eu}
	\\
\end{center}

\par\vspace{-10pt}
\rule{\linewidth}{0.4pt}
\par\vspace{7pt}

{\large\bfseries À propos}
\par\vspace{-8pt}
\rule{\linewidth}{0.4pt}
\par\vspace{5pt}

Ingénieur logiciel orienté systèmes, spécialisé en C/C++ (optimisation de performance, programmation réseau, Linux) et Java, avec un intérêt marqué pour la {\bfseries finance de marché} et les systèmes de trading.
\\
Formé à la comptabilité, à l'économie managériale et aux mathématiques à McGill University (Montréal).
\\
Contributeur open-source (plus de 10 000 lignes de code sur des packages critiques). En préparation de l'agrégation d'informatique {\bfseries (algorithmique des graphes, complexité, automates)}.
\par\vspace{5pt}

{\large\bfseries Compétences}
\par\vspace{-8pt}
\rule{\linewidth}{0.4pt}
\par\vspace{5pt}

\hspace*{5pt}{\bfseries Maîtrise}: C/C++, Java, Golang, Python, SQL (PostgreSQL), Lua, Shell
\\
\hspace*{5pt}{\bfseries Bonne connaissance}: Assembleur, Haskell, OCaml
\\
\hspace*{5pt}{\bfseries Outils}: gdb, Ghidra, Valgrind, Git, Docker, Kubernetes, CI/CD, Environnement Linux
\\
\hspace*{5pt}{\bfseries Langues}: Français, Anglais (C1, professionnel)

\par\vspace{5pt}

{\large\bfseries Expérience Professionnelle}
\par\vspace{-8pt}
\rule{\linewidth}{0.4pt}
\par\vspace{5pt}

{\bfseries Gno.land - Ingénieur Logiciel Junior (Open-Source)}
\begin{itemize}[noitemsep, topsep=0pt]
	\item Développement de 6+ packages core en Gnolang (variante de Go pour smart contracts), dont une compatibilité EVM (traduction d'opcodes en instructions Go).
	\item Détection et correction de {\bfseries 6 bugs critiques} causant des crashs runtime ; {\bfseries ~10k lignes de code} contribuées.
	\item Rédaction de tests exhaustifs sous des politiques strictes de CI/CD et de revue de code.
\end{itemize}

\par\vspace{5pt}

{\bfseries Faurecia / Forvia - Ingénieur Démo Logicielle}
\begin{itemize}[noitemsep, topsep=0pt]
	\item Conception d'un {\bfseries moteur de calcul d'itinéraire GPS} et d'une interface de recherche de POI, intégrés à une application démo présentée au {\bfseries CES 2024}.
	\item Correction d'un ancien code {\bfseries C++} pour le faire communiquer avec les bons drivers matériels (LED) ; scripts Python de pilotage de hardware.
	\item Livraison de fonctionnalités prêtes pour la production {\bfseries en avance sur le planning}, salué pour la rapidité et la qualité du code.
\end{itemize}

\par\vspace{5pt}

{\bfseries Epitech / IONIS - Assistant Technique \& Formateur}
\begin{itemize}[noitemsep, topsep=0pt]
	\item Enseignement à {\bfseries plus de 400 étudiants} en C, C++, asm et Haskell ; correction et évaluation des rendus.
	\item Création de {\bfseries contenus pédagogiques interactifs} (quiz, cours inversés, tutoriels) et contribution à un {\bfseries taux de réussite historiquement élevé} au sein des promotions coachées par l'ASTEK.
\end{itemize}

\par\vspace{5pt}

{\large\bfseries Projets Personnels}
\par\vspace{-8pt}
\rule{\linewidth}{0.4pt}
\par\vspace{5pt}

\underline{\bfseries \href{https://github.com/paulogarithm/cclass}{cclass} -- Moteur de mutation mémoire à l'exécution}
\par
\begin{adjustwidth}{5pt}{0pt}
    Framework en C pur modifiant dynamiquement des instructions assembleur compilées directement en mémoire. \\
    Utilisation d'appels système {\bfseries \texttt{mmap}} avec des flags de protection précis pour contourner le surcoût du dispatch par vtable. \\
    Application de règles d'alignement mémoire strictes pour garantir une stabilité absolue sur les cibles d'exécution matérielles.
\end{adjustwidth}
\par
{\bfseries Compétences :} C, systèmes Linux, assembleur x86, mmap, sûreté mémoire, débogage bas niveau

\par\vspace{8pt}

\underline{\bfseries \href{https://github.com/paulogarithm/blorox}{BloRox} -- Moteur 3D temps réel en C++}
\par
\begin{adjustwidth}{5pt}{0pt}
	Architecture ECS orientée données: {\bfseries 100k entités animées}, perfs de $\sim$6 ms de moteur et $\sim$6 ms de rendu par image. \\
	Alignement mémoire et extensions SIMD du processeur pour la vectorisation ; migration prévue vers {\ttfamily std::simd} (C++26). \\
	Prototypage d'une sérialisation réseau par réflexion statique C++26 (compilateur expérimental).
\end{adjustwidth}
\par
{\bfseries Compétences :} C++23/26, ECS, alignement mémoire, SIMD, optimisation de performance, conception de moteur

\par\vspace{8pt}

\underline{\bfseries Zappy -- Serveur TCP multi-clients en C}
\par
\begin{adjustwidth}{5pt}{0pt}
    Serveur TCP mono-thread à boucle événementielle ({\bfseries \texttt{select}}), codé seul : buffer borné par client et mise en quarantaine des clients qui le saturent ; aucune erreur ni fuite sous Valgrind ; tests unitaires et fonctionnels.
\end{adjustwidth}
\par
{\bfseries Compétences :} C, sockets TCP, select, gestion mémoire, Valgrind, tests unitaires et fonctionnels

\par\vspace{8pt}

\underline{\bfseries Sledge -- Moteur de migration, synchronisation et orchestration multi-cloud}
\par
\begin{adjustwidth}{5pt}{0pt}
    Outil générique de migration multi-cloud gérant des structures de données complexes et des interactions avec des bases de données. \\
    Plateforme backend complète en Go, optimisée pour des flux de données à haut volume, via une couche d'abstraction cloud personnalisée ({\ttfamily provicol}) supportant AWS S3 et GCP.
\end{adjustwidth}
\par
{\bfseries Compétences :} Go, Python, SQL, intégrité des données, cloud computing, automatisation système

\par\vspace{8pt}

{\large\bfseries Formation}
\par\vspace{-8pt}
\rule{\linewidth}{0.4pt}
\par\vspace{5pt}

{\bfseries Epitech, Paris} \hfill 2022 -- 2027 \\
Diplôme d'ingénieur en technologies de l'information (diplôme de niveau Master) \\
Certification Niveau 7 - Expert en ingénierie logicielle

\vspace{6pt}

{\bfseries McGill University, Montréal} \hfill 2025 -- 2026 \\
Certificat en Management -- GPA : 3.57/4.0 \\
Cours suivis : Comptabilité, Managerial Economics \& Analysis, Business, Mathématiques, Gestion de projet, Marketing

\end{document}
